%-------------------------------------------------------------------------------
%	SECTION TITLE
%-------------------------------------------------------------------------------
\cvsection{Опыт работы}

%-------------------------------------------------------------------------------
%	CONTENT
%-------------------------------------------------------------------------------
\begin{cventries}

%---------------------------------------------------------
  \cventry
    {Инженер-исследователь \textnormal{· основное место работы, 1.0 ставки}} % Job title
    {\href{https://github.com/brain-lab-research}{BRAIn Lab}, МФТИ} % Organisation
    {Москва, Россия} % Location
    {Сентябрь 2023 — н.в.} % Date(s)
    {
      \begin{cvitems}
        \item {Формально: лаборатории математических методов оптимизации, МФТИ-Яндекс и фундаментальных исследований ИИ}
        \item {Предобучение (pretrain) больших языковых моделей и эффективные методы оптимизации для него}
        \item {Оптимизаторы нового поколения (Muon, SignSGD), методы нулевого порядка, расписания warm-up}
        \item {AI Scientist: автономные агенты для проведения научных исследований; научные \href{https://rutube.ru/channel/16197436/videos/}{семинары}}
      \end{cvitems}
    }

%---------------------------------------------------------
  \cventry
    {Руководитель направления \textnormal{(ранее стажёр) · по совместительству, 0.5 ставки}} % Job title
    {\href{https://sberlabs.com/laboratories/sber-ai-lab}{Сбербанк} — Центр практического ИИ / AI Lab} % Organisation
    {Москва, Россия} % Location
    {Сентябрь 2024 — н.в.} % Date(s)
    {
      \begin{cvitems}
        \item {Вырос от стажёра до руководителя направления в центре практического искусственного интеллекта}
        \item {Эффективное дообучение больших языковых моделей, разработка LoRA-адаптеров (WeightLoRA, ACL 2026)}
        \item {Внедрение ИИ-решений в продуктовые задачи Банка}
      \end{cvitems}
    }

%---------------------------------------------------------
  \cventry
    {Старший лаборант \textnormal{(ранее лаборант) · по совместительству, 0.5 ставки}} % Job title
    {\href{https://www.ispras.ru}{ИСП РАН} — Исследовательский центр доверенного ИИ} % Organisation
    {Москва, Россия} % Location
    {Февраль 2024 — н.в.} % Date(s)
    {
      \begin{cvitems}
        \item {Разработка методов федеративного обучения и доверенного искусственного интеллекта}
        \item {Организация и проведение распределённого обучения больших языковых моделей}
      \end{cvitems}
    }

%---------------------------------------------------------
  \cventry
    {Research Assistant} % Job title
    {\href{https://mbzuai.ac.ae/research-department/machine-learning-department/}{MBZUAI}, ML Department} % Organisation
    {Абу-Даби, ОАЭ} % Location
    {Февраль 2025 — Апрель 2025} % Date(s)
    {
      \begin{cvitems}
        \item {Международный исследовательский опыт в одном из ведущих мировых университетов по искусственному интеллекту}
        \item {Совместная работа с международной командой (M. Tak{\'a}{\v{c}}, S. Horv{\'a}th) над эффективными методами обучения матричных параметров (Muon-like) и инвариантной оптимизацией}
        \item {По итогам сотрудничества — серия совместных статей по оптимизаторам (OPLoRA, Preconditioned Norms, LionMuon)}
      \end{cvitems}
    }

%---------------------------------------------------------
  \cventry
    {Младший научный сотрудник \textnormal{(ранее стажёр-исследователь)}} % Job title
    {\href{http://iitp.ru/ru/researchlabs/1034.htm}{ИППИ РАН} — лаб. распределённых вычислительных систем} % Organisation
    {Москва, Россия} % Location
    {Февраль 2024 — Октябрь 2025} % Date(s)
    {
      \begin{cvitems}
        \item {Исследование петель обратной связи в системах машинного обучения}
        \item {Распределённые и марковские методы оптимизации}
      \end{cvitems}
    }

%---------------------------------------------------------
  \cventry
    {Стажёр-исследователь} % Job title
    {\href{https://new.skoltech.ru}{Сколтех}} % Organisation
    {Москва, Россия} % Location
    {Июнь 2022 — Сентябрь 2022} % Date(s)
    {
      \begin{cvitems}
        \item {Методы детектирования разладок временных рядов с помощью машинного обучения}
      \end{cvitems}
    }

\end{cventries}
